% GENERATED FILE. Edit canonical lessons, then run npm run generate:books.
% canonical-source-hash: fnv1a64:52136571f4872290
% canonical-lessons: LA-C460-turba, LA-C460-frequens, LA-C460-frequento, LA-C460-commeo, LA-C460-suburbanus

\chapter{Crowd, Crowded, To throng, To go back and forth, Near the city}
\label{ch:la-crowd-crowded-to-throng-to-go-back-and-forth-near-the-city}
\hlchaptermodality{460}

\begin{chapteropening}
\textbf{By the end of this chapter:} \emph{I can describe a busy town: a crowd, crowded streets, people thronging and going back and forth, and places near the city.}

Complete the last lesson of chapter 460: I can describe a busy town: a crowd, crowded streets, people thronging and going back and forth, and places near the city.
\end{chapteropening}

\section[turba, turbae]{turba, turbae — a crowd, a throng; an uproar}

\label{lesson:LA-C460-turba}



\begin{quote}
Before the new one: say the Latin for the sacred boundary, then the Latin for to build up.
\end{quote}

\subsection*{You'll want to know: turba, turbae}
\textbf{turba, turbae} — \textquotedblleft{}a crowd, a throng; an uproar\textquotedblright{}.

\begin{quote}
In viā turba magna est, et raeda lentē it.
\end{quote}

\begin{quote}
\emph{There is a big crowd in the street, and the carriage goes slowly.}
\end{quote}

First declension, feminine: \emph{turba, turbae f.} First disorder and noise, then the crowd that makes it. English \emph{disturb} and \emph{turbulent} belong to the same family.

\begin{cousinweb}
\emph{turba} is usually linked with Greek \emph{tyrbē} 'turmoil, disorder' (EtymDB; Lewis \& Short); whether it was borrowed or is a cousin is not settled. English \emph{disturb} and \emph{turbulent} come from words built on it.
\end{cousinweb}

\begin{grammarlens}[title={What you've built}]
One more word for telling, explaining and describing at B1.
\end{grammarlens}

\subsection*{Your turn}
\begin{itemize}
\raggedright
  \item \emph{Say it:} \emph{turba, turbae}
  \item \emph{Say it:} \emph{turba, turbae}, once more
  \item \emph{Say it:} a sentence of your own with \emph{turba, turbae}
  \item \emph{Recall:} say the Latin for the sacred boundary, then the Latin for to build up, then say \emph{turba, turbae} again
\end{itemize}

\begin{tcolorbox}[breakable,colback=teal!4,colframe=teal!35!black,title={Before you move on}]
What does turba, turbae mean? (\textquotedblleft{}A crowd, a throng; an uproar\textquotedblright{}.) Say it once more.
\end{tcolorbox}
\section[frequēns, frequentis]{frequēns, frequentis — crowded, full of people; frequent}

\label{lesson:LA-C460-frequens}



\begin{quote}
Before the new one: say the Latin for to build up, then the Latin for a crowd.
\end{quote}

\subsection*{You'll want to know: frequēns, frequentis}
\textbf{frequēns, frequentis} — \textquotedblleft{}crowded, full of people; frequent\textquotedblright{}.

\begin{quote}
Forum hodiē frequēns est, quod mercātus est.
\end{quote}

\begin{quote}
\emph{The forum is crowded today, because it is market day.}
\end{quote}

An adjective of one ending: \emph{frequēns}, genitive \emph{frequentis}. First 'packed close', then 'often'. English \emph{frequent} comes from it.

\begin{cousinweb}
\emph{frequēns} is usually traced to PIE \emph{b\textsuperscript{h}rek-} 'to cram together' (EtymDB; Lewis \& Short compare \emph{farciō} 'to stuff'). English \emph{frequent} comes from it.
\end{cousinweb}

\begin{grammarlens}[title={What you've built}]
One more word for telling, explaining and describing at B1.
\end{grammarlens}

\subsection*{Your turn}
\begin{itemize}
\raggedright
  \item \emph{Say it:} \emph{frequēns, frequentis}
  \item \emph{Say it:} \emph{frequēns, frequentis}, once more
  \item \emph{Say it:} a sentence of your own with \emph{frequēns, frequentis}
  \item \emph{Recall:} say the Latin for to build up, then the Latin for a crowd, then say \emph{frequēns, frequentis} again
\end{itemize}

\begin{tcolorbox}[breakable,colback=teal!4,colframe=teal!35!black,title={Before you move on}]
What does frequēns, frequentis mean? (\textquotedblleft{}Crowded, full of people; frequent\textquotedblright{}.) Say it once more.
\end{tcolorbox}
\section[frequentō, frequentāre]{frequentō, frequentāre — to throng, to crowd into; to visit often}

\label{lesson:LA-C460-frequento}



\begin{quote}
Before the new one: say the Latin for a crowd, then the Latin for crowded.
\end{quote}

\subsection*{You'll want to know: frequentō, frequentāre}
\textbf{frequentō, frequentāre} — \textquotedblleft{}to throng, to crowd into; to visit often\textquotedblright{}.

\begin{quote}
Cīvēs thermās post merīdiem frequentant.
\end{quote}

\begin{quote}
\emph{The citizens throng the baths in the afternoon.}
\end{quote}

\emph{frequentō, frequentāre, frequentāvī, frequentātum}, with an accusative. English \emph{frequent} as a verb comes from it.

\begin{cousinweb}
\emph{frequentō} is built on \emph{frequēns, frequentis} 'crowded, frequent' (EtymDB; Lewis \& Short).
\end{cousinweb}

\begin{grammarlens}[title={What you've built}]
One more word for telling, explaining and describing at B1.
\end{grammarlens}

\subsection*{Your turn}
\begin{itemize}
\raggedright
  \item \emph{Say it:} \emph{frequentō, frequentāre}
  \item \emph{Say it:} \emph{frequentō, frequentāre}, once more
  \item \emph{Say it:} a sentence of your own with \emph{frequentō, frequentāre}
  \item \emph{Recall:} say the Latin for a crowd, then the Latin for crowded, then say \emph{frequentō, frequentāre} again
\end{itemize}

\begin{tcolorbox}[breakable,colback=teal!4,colframe=teal!35!black,title={Before you move on}]
What does frequentō, frequentāre mean? (\textquotedblleft{}To throng, to crowd into; to visit often\textquotedblright{}.) Say it once more.
\end{tcolorbox}
\section[commeō, commeāre]{commeō, commeāre — to go back and forth, to come and go}

\label{lesson:LA-C460-commeo}



\begin{quote}
Before the new one: say the Latin for crowded, then the Latin for to throng.
\end{quote}

\subsection*{You'll want to know: commeō, commeāre}
\textbf{commeō, commeāre} — \textquotedblleft{}to go back and forth, to come and go\textquotedblright{}.

\begin{quote}
Mercātōrēs inter portum et forum commeant.
\end{quote}

\begin{quote}
\emph{Merchants go back and forth between the harbour and the forum.}
\end{quote}

\emph{commeō, commeāre, commeāvī, commeātum}. The noun \emph{commeātus} is the coming and going, then supplies and leave.

\begin{cousinweb}
\emph{commeō} is built from \emph{con-} and \emph{meō} 'to go, to pass' (EtymDB; Lewis \& Short).
\end{cousinweb}

\begin{grammarlens}[title={What you've built}]
One more word for telling, explaining and describing at B1.
\end{grammarlens}

\subsection*{Your turn}
\begin{itemize}
\raggedright
  \item \emph{Say it:} \emph{commeō, commeāre}
  \item \emph{Say it:} \emph{commeō, commeāre}, once more
  \item \emph{Say it:} a sentence of your own with \emph{commeō, commeāre}
  \item \emph{Recall:} say the Latin for crowded, then the Latin for to throng, then say \emph{commeō, commeāre} again
\end{itemize}

\begin{tcolorbox}[breakable,colback=teal!4,colframe=teal!35!black,title={Before you move on}]
What does commeō, commeāre mean? (\textquotedblleft{}To go back and forth, to come and go\textquotedblright{}.) Say it once more.
\end{tcolorbox}
\section[suburbānus, suburbāna, suburbānum]{suburbānus, suburbāna, suburbānum — near the city, suburban}

\label{lesson:LA-C460-suburbanus}



\begin{quote}
Before the new one: say the Latin for to throng, then the Latin for to go back and forth.
\end{quote}

\subsection*{You'll want to know: suburbānus, suburbāna, suburbānum}
\textbf{suburbānus, suburbāna, suburbānum} — \textquotedblleft{}near the city, suburban\textquotedblright{}.

\begin{quote}
Dīves vīllam suburbānam habet et saepe ibi cēnat.
\end{quote}

\begin{quote}
\emph{The rich man has a villa near the city and often dines there.}
\end{quote}

An adjective of three endings, from \emph{sub} 'near, under' and \emph{urbs} 'city'. \emph{Suburbānum}, a villa near the city. English \emph{suburban} comes from it.

\begin{cousinweb}
\emph{suburbānus} is built from \emph{sub} 'under, close to' and \emph{urbānus} 'of the city', from \emph{urbs} (Lewis \& Short). English \emph{suburban} comes from it.
\end{cousinweb}

\begin{grammarlens}[title={What you've built}]
One more word for telling, explaining and describing at B1.
\end{grammarlens}

\subsection*{Your turn}
\begin{itemize}
\raggedright
  \item \emph{Say it:} \emph{suburbānus, suburbāna, suburbānum}
  \item \emph{Say it:} \emph{suburbānus, suburbāna, suburbānum}, once more
  \item \emph{Say it:} a sentence of your own with \emph{suburbānus, suburbāna, suburbānum}
  \item \emph{Recall:} say the Latin for to throng, then the Latin for to go back and forth, then say \emph{suburbānus, suburbāna, suburbānum} again
\end{itemize}

\begin{tcolorbox}[breakable,colback=teal!4,colframe=teal!35!black,title={Before you move on}]
What does suburbānus, suburbāna, suburbānum mean? (\textquotedblleft{}Near the city, suburban\textquotedblright{}.) Say it once more.
\end{tcolorbox}
